\clearpage
\fancyhead[L]{\small DE LATITUDINE SPECULAE MANHEMIENSIS.}
\hypertarget{native-602-lower}{}
\begin{center}\small Göttingische gelehrte Anzeigen. No. 76. Pages 753..756. 12 May 1817.\\[3pt]\rule{.72\linewidth}{.3pt}\end{center}

\emph{De latitudine speculae Manhemiensis autore} H.\,C.\ \textsc{Schumacher},
\emph{astronomiae professore et regiae societatis scientiarum Havniensis
socio.} Copenhagen, 1816. Printed by J. F. Schulz. 56 pages in large quarto.

\begingroup\brokenpenalty=10000\relax
For measuring zenith distances, the Mannheim Observatory has an eight-foot mural \SourceLeaf{603}quadrant by \textsc{Bird}, a nine-foot zenith sector by \textsc{Sisson}, and since 1811 a three-foot repeating circle by \textsc{Reichenbach} as well. No arrangements have yet been made to mount the last instrument. The author, who was director of the observatory from 1813 to 1815, was therefore restricted to the first two instruments in determining the observatory’s latitude as accurately as the available means permitted. The zenith sector was used to determine the collimation error of the mural quadrant. Since the mural quadrant faces south and cannot be remounted in the opposite orientation, an independent determination of the latitude from fixed stars was impossible. The only course left was to take as a basis the declinations of the fixed stars observed, according to the most reliable and recent determinations by other astronomers. One hundred and fifty-two observations made with the zenith sector from 27 January 1814 to 4 March 1815 served to determine the mean zenith distances for eighteen stars, with the distances reduced to the beginning of 1815. These, however, required a correction, since an examination showed that the degree divisions on the zenith sector were too small and that the measured zenith distances were consequently too large.
To determine the law governing these corrections, the author used a beam compass to measure the chords of several arcs on the limb of the sector and at the same time the distance from the centre of motion.
The formula for the correction that Mr. \textsc{Bessel} in Königsberg derived from these data is
\[
1''3892\,z
\]
where $z$ denotes the zenith distance in degrees. It represents those measurements fairly well, and the small differences may reasonably be regarded as measurement errors.
The author has meanwhile preferred the formula
\[
0''4225\left(z+\tfrac{1}{10}zz+\tfrac{1}{20}z^{3}\right)
\]
which shows still better agreement.
This assumes, then, that the degree divisions on the zenith sector are too small and vary in size.
Strictly speaking, however, the second term is inadmissible here because the zero point bisects all the arcs whose chords were measured.
The formula
\[
0''585\left(z+\tfrac{1}{28}z^{3}\right)
\]
would also have given just as good agreement.
As for the final result, it makes almost no difference which of these correction formulas one uses. The zenith distances found in this way were then used to determine the collimation error of the mural quadrant, with which the same stars had been observed. Using this error, the zenith distances of 36 other stars observed with the mural quadrant were calculated. Comparison with the declinations determined by \textsc{Piazzi}, \textsc{Pond} and \textsc{Oriani} gave $49^{\circ}\,29'\,13''5$ for the latitude of the Mannheim Observatory as the result of 284 observations. Of the 18 stars observed with the zenith sector itself, the author could use only nine in the same way, since he lacked equally reliable determinations of the declinations of the others. From 87 observations, these nine gave him a latitude of $49^{\circ}\,29'\,14''2$. The reviewer reduced the observations of the other nine using the declinations in \textsc{Piazzi}'s new catalogue, which the author had not yet been able to use, and found the following results:
\begin{center}
\fontsize{9}{11}\selectfont
\setlength{\tabcolsep}{3pt}
\begin{tabular}{l|c|c||l|c|c}
 & Latitude & \begin{tabular}{c}Number of\\observations\end{tabular} & & Latitude & \begin{tabular}{c}Number of\\observations\end{tabular} \\
\hline
$\sigma$ Perseus & $49^{\circ}\,29'\,10''4$ & 8 & 725 Perseus & $49^{\circ}\,29'\,11''5$ & 9 \\
$\psi$ Perseus & 8.3 & 9 & 7 Perseus & 8.2 & 8 \\
$\chi$ Ursa Major & 8 3 & 5 & $F$ Ursa Major & 9.8 & 5 \\
$\theta$ Perseus & 15.2 & 5 & $\gamma$ Perseus & 8.2 & 13 \\
$\iota$ Ursa Major & 8.0 & 5 & \multicolumn{3}{|l}{} \\
\end{tabular}
\end{center}
\noindent\SourceLeaf{604}For the star the author designated as 725 Perseus, whose designation could not be found, entry II, 253 in \textsc{Piazzi}'s catalogue was adopted. These 65 observations thus give a mean of $49^\circ\,29'\,9''8$. It follows that all 152 observations with the zenith sector give a mean of $49^\circ\,29'\,12''3$, and all 436 observations of fixed stars give a mean of $49^\circ\,29'\,13''1$. The 66 observations of the Sun made by the author in 1814 give a mean only a few tenths of a second greater, but for certain local reasons they cannot lay claim to the same reliability. This result from the author's observations will now have to be regarded as the most certain until \textsc{Reichenbach}'s circle has been mounted and used to make the new latitude determination from circumpolar stars, the only one really worthy of such an observatory\enspace Furthermore, the above result does agree exactly with \textsc{Barry}'s last reported value, but since that astronomer had not made known how he obtained it, so that its degree of accuracy could not be judged, Mr. \textsc{Schumacher} deserves the thanks of all astronomers for publishing his observations in full.
\par\clearpage\endgroup
